\chapter{Bat Bites}

\section*{Definition and Overview}
Bat bites are puncture wounds or scratches caused by the teeth or claws of a bat. Although bats are small, their teeth are sharp, and even a tiny bite or scratch may not bleed or be noticed. The main medical concern is not the wound itself but the risk of infection, particularly rabies, which is a nearly always fatal viral disease of the nervous system. Because of this, any contact with a bat that may have broken the skin requires prompt medical assessment, even if the wound looks trivial.

\section*{Epidemiology}
Bat bites are uncommon compared with bites from dogs and cats, but they are a recognised risk for people who handle bats, such as wildlife carers, vets, and people who find sick or injured bats. In Australia, bats (flying foxes and microbats) can carry the Australian bat lyssavirus, a close relative of rabies, and a small number of human cases have been reported. Worldwide, bats are a major reservoir of rabies virus. Because bites may go unnoticed (for example, while sleeping), people may not realise they have been exposed.

\section*{Aetiology and Pathophysiology}
\begin{itemize}
    \item \textbf{Bite mechanism:} a bat's small, sharp teeth can produce puncture wounds; scratches from claws can also break the skin
    \item \textbf{Saliva exposure:} the virus is present in bat saliva and enters the body through broken skin, the eyes, nose, or mouth
    \item \textbf{Australian bat lyssavirus (ABLV):} the relevant virus in Australia, closely related to rabies virus
    \item \textbf{Rabies virus:} found in bats in many other countries, including the Americas, Asia, and Africa
    \item \textbf{Mechanism of disease:} the virus travels along nerves from the bite site to the brain, causing fatal encephalitis; the incubation period is usually weeks to months
\end{itemize}

\section*{Clinical Features}
\begin{itemize}
    \item A tiny puncture wound, scratch, or graze that may be barely visible or may bleed very little
    \item Local pain, redness, or swelling at the wound site
    \item The bite may be entirely unnoticed, especially after sleeping or handling a bat
    \item If rabies or ABLV develops (rare, if untreated): fever, headache, tingling at the wound site, confusion, agitation, paralysis, and difficulty swallowing
\end{itemize}

\section*{Differential Diagnosis}
\begin{itemize}
    \item Bites or scratches from other animals (dogs, cats, rodents) that also carry infection risk
    \item Other puncture wounds from sharp objects that may be infected
    \item Tick or insect bites that cause a small skin mark
    \item If rabies is suspected, other causes of encephalitis, such as other viruses, meningitis, or tetanus, must be considered
\end{itemize}

\section*{When to Seek Medical Care}
Any bite, scratch, or exposure to bat saliva warrants urgent medical assessment, even if the wound is tiny or the bat appears healthy. Wash the wound immediately with soap and water for at least five minutes, then seek care the same day.

\textbf{Call 000 (or local emergency services) for:}
\begin{itemize}
    \item Any bite or scratch from a bat, or contact with bat saliva that may have entered the eyes, nose, or mouth
    \item A bat found in a room with a sleeping child, an elderly person, or a person with impaired awareness
    \item Symptoms of nervous system illness, such as confusion, weakness, or difficulty swallowing, after bat exposure
    \item Signs of a serious wound infection: spreading redness, severe pain, or fever
\end{itemize}

\section*{Diagnosis and Investigations}
\begin{itemize}
    \item \textbf{History:} details of the bat contact, whether the wound bled, and whether the bat was handled or found in a room with a sleeping person
    \item \textbf{Examination:} inspection of the wound and assessment for any breaks in the skin
    \item \textbf{Risk assessment:} each exposure is assessed for the need for rabies post-exposure prophylaxis (vaccine and immunoglobulin)
    \item \textbf{Bat testing (where possible):} the bat may be tested for lyssavirus to guide treatment, but prophylaxis should not be delayed while waiting
    \item \textbf{Wound care:} cleaning and assessment of the wound; tetanus vaccination is updated if needed
    \item \textbf{No routine blood tests:} diagnosis of lyssavirus infection in a person is made clinically and by specialised laboratory testing if symptoms develop
\end{itemize}

\section*{Management}
\begin{itemize}
    \item \textbf{Immediate first aid:} wash the wound thoroughly with soap and water for at least five minutes; apply antiseptic
    \item \textbf{Rabies post-exposure prophylaxis:} a course of rabies vaccine, often combined with rabies immunoglobulin injected into and around the wound, given as soon as possible after exposure
    \item \textbf{Tetanus prophylaxis:} update tetanus vaccination if not current
    \item \textbf{Wound care:} dressing the wound; antibiotics are only used if there are signs of bacterial infection
    \item \textbf{Observation:} monitoring over the following months for any neurological symptoms
    \item \textbf{Public health notification:} bat exposures are reported to public health authorities to coordinate care
\end{itemize}

\section*{Complications}
\begin{itemize}
    \item Rabies and Australian bat lyssavirus encephalitis -- nearly always fatal if symptoms develop; prevention before symptoms is essential
    \item Bacterial wound infection, including cellulitis or, rarely, deeper infection (for example, in the joints or bone of the hand)
    \item Tetanus if vaccination is not up to date
    \item Anxiety and distress about the exposure and its potential consequences
\end{itemize}

\section*{Prognosis and Outcomes}
The outcome is excellent when treatment is given promptly after exposure: modern post-exposure vaccination is highly effective, and people who complete it do not develop the disease. Once symptoms of lyssavirus infection begin, the disease is almost always fatal, which is why urgent treatment after any possible exposure is so important. Healing of the small wound itself is usually uneventful.

\section*{Prevention and Self-Care}
\begin{itemize}
    \item Never handle bats with bare hands; use thick gloves if handling is necessary (best left to trained professionals)
    \item Do not touch sick, injured, or dead bats
    \item Keep bats out of living areas; exclude them from roofs and verandas, and seal gaps where they can enter
    \item If a bat is found in a room where someone was sleeping, seek medical advice even if no bite is visible
    \item Teach children not to touch bats and to report any contact
    \item People at high risk (bat handlers, vets, wildlife workers) may be offered pre-exposure rabies vaccination
    \item Wash any wound from a bat bite immediately with soap and water, then seek medical attention the same day
\end{itemize}

\section*{Sources and Further Reading}
The following sources provide reliable, up-to-date information on bat bites and rabies prevention.
\begin{itemize}
    \item Centers for Disease Control and Prevention. \textit{Rabies.} https://www.cdc.gov
    \item World Health Organization. \textit{Rabies.} https://www.who.int
    \item NHS. \textit{Rabies.} https://www.nhs.uk/conditions/
    \item MSD Manual. \textit{Rabies.} https://www.msdmanuals.com
    \item MedlinePlus. \textit{Rabies.} https://medlineplus.gov
\end{itemize}
